\scanpage{54}

Du fait que $i^!$ est l'adjoint à droite
\margin{\deleted{\unsure{commutant} aux \unsure{produits}}}
du foncteur exact \deleted{à gauche} $i_!$, on
déduit facilement la \uline{formule de}
\uline{dualité} \uline{pour $i$}
\[
\boxed{\operatorname{Hom}(G^\bullet,\mathbb{R}i^!(F^\bullet))
\simeq\operatorname{Hom}(\mathbb{R}_!i_*\unsure{G^\bullet|U^\bullet}),F^\bullet)}
\]
\margin{$G$}
(isom bifonctoriel).

\[
\begin{aligned}
\operatorname{Hom}(G^\bullet,\mathbb{R}i^!(F^\bullet))
&=\operatorname{Hom}(\unsure{C^\bullet|G^\bullet},i^!(C^\bullet(F^\bullet)))\\
&=H^0(\operatorname{Hom}^{\bullet}(G^\bullet,i^!(C^\bullet(F^\bullet)))
=H^0(\operatorname{Hom}(i_!(G^\bullet),C^\bullet(F^\bullet)))\\
&\deleted{=\operatorname{Hom}\unsure{(G^\bullet(G))}}\\
&=\operatorname{Hom}(\mathbb{R}i_!(G^\bullet),C^\bullet(F^\bullet))
\end{aligned}
\]

Quand on \unsure{explicite} la flèche $\unsure{-|\longrightarrow}$, on trouve
qu'elle est associée à un hom can.
\[
\operatorname{Tr}_i(F^\bullet):\quad
\mathbb{R}_!i_*\deleted{\unsure{G}}\inserted{(\mathbb{R}i^!(F^\bullet))}
\longrightarrow F^\bullet
\]
de la \unsure{même} façon que dans le cas du th.
de dualité globale pour un morphisme
lisse.
\mspage{53}
